\section{Discussion}
\label{sec:discussion}

\paragraph{Behavioral Suppression vs.\ Representational Erasure.}
Our central finding is a consistent gap between behavioral and representational
evidence of forgetting.
All eight unlearned models show substantial drops in generation and logit-based
accuracy on WMDP-bio, suggesting successful behavioral suppression.
However, probes trained on the corresponding hidden states continue to classify
the forget domain with high accuracy.
This gap indicates that current gradient-based unlearning methods do not erase
knowledge representations---they primarily modify the model's \emph{output mapping}
from those representations to generated text.
Concretely, the model still encodes ``which answer is biologically correct'' in
mid-to-late layer activations, but suppresses this information before it reaches
the output layer.

\paragraph{Method Differences.}
Not all methods are equally effective at behavioral suppression, nor do they
differ uniformly in their representational impact.
Some methods (e.g., RMU, RMU-LAT) are specifically designed to disrupt internal
representations by steering them toward noise, and may therefore show slightly
lower probe accuracy than methods based purely on output-level gradient differences
(e.g., GradDiff).
Full rankings are reported in \cref{tab:probes-bio}.
Even the most representation-targeted method, however, does not close the gap to
random chance for probes, suggesting that current representation-disruption
objectives are insufficient to fully erase knowledge.

\paragraph{Cross-Probe Transfer.}
The high cross-probe accuracy in both directions (\emph{Base $\to$ Method} and
\emph{Method $\to$ Base}) provides strong evidence that unlearning does not
substantially alter the linear geometry of knowledge-encoding directions.
A probe trained entirely on base model hidden states transfers to unlearned model
hidden states with minimal accuracy drop, meaning the same linear subspace that
discriminates biosecurity knowledge in the base model remains discriminative after
unlearning.
This is consistent with the view that gradient-based unlearning fine-tunes the
logit-generation pathway (later MLP weights and the unembedding matrix) while
leaving earlier representational layers relatively intact.

\paragraph{Domain Generalization.}
The WMDP-cyber results reveal that representational preservation is not confined to
the explicit unlearning target.
Cybersecurity knowledge, which is not a target of any of the eight methods, shows a
similar pattern: generation accuracy is largely preserved (as expected), and so is
probe accuracy.
This confirms that the phenomenon is not an artifact of the specific forget domain.

\paragraph{Implications for Safety Evaluation.}
Current safety evaluations of unlearned LLMs rely almost entirely on behavioral
benchmarks such as WMDP accuracy.
Our results demonstrate that such evaluations can be misleading: a model may appear
to have forgotten on behavioral metrics while retaining the knowledge in its internal
representations.
An adversary with white-box access (e.g., a researcher testing a shared model
checkpoint, or an insider with API access to activations) could potentially build
a probe on a small set of labeled examples and use it to extract suppressed knowledge.
This motivates the development of \emph{representational} unlearning methods and
corresponding representational auditing protocols as a complement to behavioral
benchmarks.

\paragraph{Limitations.}
\begin{itemize}[leftmargin=1.5em, itemsep=2pt]
    \item \emph{Model family.} We study a single model family (Llama-3-8B-Instruct)
    and a single set of unlearning checkpoints.
    Whether the findings generalize to other architectures (e.g., Mistral, Gemma)
    or unlearning implementations remains open.

    \item \emph{Probe complexity.} Our probes are shallow (logistic regression,
    random forest, AdaBoost).
    The fact that such simple probes already achieve high accuracy suggests the
    knowledge is linearly decodable; more complex probes could only do better.

    \item \emph{True/False framing.} We evaluate on a claim-verification
    reformulation of WMDP rather than the original multiple-choice format.
    Generation accuracy on this framing may differ from the standard MCQ
    accuracy reported in the original WMDP leaderboard.

    \item \emph{No adaptive adversary.} We consider a non-adaptive probe---one
    trained on a fixed small labeled set.
    An adaptive adversary with access to more resources or knowledge of the model
    internals might achieve even higher accuracy.
\end{itemize}

\paragraph{Future Work.}
Future work should develop unlearning objectives that explicitly minimize probe
accuracy (not just generation accuracy), e.g., by adversarially training a probe
simultaneously with the unlearning optimizer.
Mechanistic studies of \emph{which} weight structures encode the forget-domain
knowledge---using causal tracing~\citep{meng2022locating} applied to the unlearning
setting---could guide more targeted erasure.
Finally, representational auditing protocols analogous to behavioral benchmarks
would provide a more complete picture of unlearning effectiveness.
